% Theorem System
% The following blocks are provided:
%   Definition:     \defn
%   Theorem:        \thm
%   Lemma:          \lem
%   Corollary:      \cor
%   Proposition:    \prop
%   Claim:          \clm
%   Fact:           \fact
%   Assumption:     \assumption
%   Proof:          \pf
%   Example:        \ex
%   Remark:         \rmk (sentence), \rmkb (block)
%
% Toggle:
%   \EnablePlainAmsthmTheoremtrue  -> native amsthm theorem environments
%   \EnablePlainAmsthmTheoremfalse -> styled tcolorbox theorem environments
%
% Label behavior for theorem commands:
%   - empty label: no \label written
%   - already prefixed (e.g. thm:foo for \thm): keep as-is
%   - otherwise auto-prefix (e.g. foo -> thm:foo)

\makeatletter
\@ifundefined{ifEnablePlainAmsthmTheorem}{
    \newif\ifEnablePlainAmsthmTheorem
    \EnablePlainAmsthmTheoremfalse
}{}

\newif\ifThemeTheoremNumberingNone
\ThemeTheoremNumberingNonefalse
\newif\ifThemeTheoremCounterWithinChapter
\ThemeTheoremCounterWithinChapterfalse
\def\ThemeTheoremCounterWithin{section}

\def\theme@theorem@chapter{chapter}
\def\theme@theorem@section{section}
\def\theme@theorem@none{none}
\@ifundefined{ThemeResolvedTheoremNumbering}{}{
    \ifx\ThemeResolvedTheoremNumbering\theme@theorem@chapter
        \@ifundefined{chapter}{
            \def\ThemeTheoremCounterWithin{section}
            \ThemeTheoremCounterWithinChapterfalse
        }{
            \def\ThemeTheoremCounterWithin{chapter}
            \ThemeTheoremCounterWithinChaptertrue
        }
    \else
        \ifx\ThemeResolvedTheoremNumbering\theme@theorem@none
            \ThemeTheoremNumberingNonetrue
            \ThemeTheoremCounterWithinChapterfalse
        \else
            \def\ThemeTheoremCounterWithin{section}
            \ThemeTheoremCounterWithinChapterfalse
        \fi
    \fi
}

\newcommand{\ThemeOptionalTheoremTitle}[1]{%
    \if\relax\detokenize{#1}\relax
    \else
        [#1]%
    \fi
}
\makeatother

\ExplSyntaxOn
\tl_new:N \l_theme_tcb_label_tl

\cs_new_protected:Npn \theme_emit_prefixed_label:nn #1#2
{
    \tl_if_blank:nF {#2}
    {
        \regex_match:nnTF { \A#1: } {#2}
        {
            \label{#2}
        }
        {
            \label{#1:#2}
        }
    }
}

\cs_new_protected:Npn \theme_normalize_tcb_label:nn #1#2
{
    \tl_clear:N \l_theme_tcb_label_tl
    \tl_if_blank:nF {#2}
    {
        \tl_set:Nn \l_theme_tcb_label_tl {#2}
        \regex_match:nnT { \A#1: } {#2}
        {
            \regex_replace_once:nnN { \A#1: } {} \l_theme_tcb_label_tl
        }
    }
}

\cs_new_protected:Npn \theme_begin_tcb_theorem:nnn #1#2#3
{
    \begin{#1}{#2}{#3}
}

\cs_new_protected:Npn \theme_run_tcb_theorem:nnnnn #1#2#3#4#5
{
    \theme_normalize_tcb_label:nn {#2} {#4}
    \exp_args:Nne \theme_begin_tcb_theorem:nnn {#1} {#3} { \tl_use:N \l_theme_tcb_label_tl }
        #5
    \end{#1}
}

\cs_new_protected:Npn \ThemeEmitPrefixedLabel #1#2
{
    \theme_emit_prefixed_label:nn {#1} {#2}
}

\cs_new_protected:Npn \ThemeRunTcbTheorem #1#2#3#4#5
{
    \theme_run_tcb_theorem:nnnnn {#1} {#2} {#3} {#4} {#5}
}
\ExplSyntaxOff

% -------------------- Native amsthm theorem family --------------------
% These environments are always available for imported paper-style LaTeX:
% \begin{theorem}, \begin{lemma}, \begin{proposition}, \begin{corollary},
% \begin{claim}, and \begin{remark} stay plain even when shortcut commands
% such as \thm use styled tcolorbox blocks.
\ifThemeTheoremNumberingNone
    \theoremstyle{definition}
    \newtheorem{definition}{Definition}
    \theoremstyle{plain}
    \newtheorem{theorem}{Theorem}
    \newtheorem{lemma}{Lemma}
    \newtheorem{proposition}{Proposition}
    \newtheorem{corollary}{Corollary}
    \newtheorem{claim}{Claim}
    \theoremstyle{remark}
    \newtheorem{remark}{Remark}
\else
    \ifThemeTheoremCounterWithinChapter
        \theoremstyle{definition}
        \newtheorem{definition}{Definition}[chapter]
        \theoremstyle{plain}
        \newtheorem{theorem}{Theorem}[chapter]
        \newtheorem{lemma}{Lemma}[chapter]
        \newtheorem{proposition}{Proposition}[chapter]
        \newtheorem{corollary}{Corollary}[chapter]
        \newtheorem{claim}{Claim}[chapter]
        \theoremstyle{remark}
        \newtheorem{remark}{Remark}[chapter]
    \else
        \theoremstyle{definition}
        \newtheorem{definition}{Definition}[section]
        \theoremstyle{plain}
        \newtheorem{theorem}{Theorem}[section]
        \newtheorem{lemma}{Lemma}[section]
        \newtheorem{proposition}{Proposition}[section]
        \newtheorem{corollary}{Corollary}[section]
        \newtheorem{claim}{Claim}[section]
        \theoremstyle{remark}
        \newtheorem{remark}{Remark}[section]
    \fi
\fi

\ifEnablePlainAmsthmTheorem
    \ifThemeTheoremNumberingNone
        \theoremstyle{plain}
        \newtheorem{fact}{Fact}
        \newtheorem{assumption}{Assumption}
    \else
        \ifThemeTheoremCounterWithinChapter
            \theoremstyle{plain}
            \newtheorem{fact}{Fact}[chapter]
            \newtheorem{assumption}{Assumption}[chapter]
        \else
            \theoremstyle{plain}
            \newtheorem{fact}{Fact}[section]
            \newtheorem{assumption}{Assumption}[section]
        \fi
    \fi
    \let\ThemeFactEnvBegin\fact
    \let\ThemeAssumptionEnvBegin\assumption
\else
    % -------------------- Styled tcolorbox theorem family --------------------
    \ifThemeTheoremNumberingNone
        \newtcbtheorem{mydefinition}{Definition}
        {
            theorem/variant={definition-body-bg}{definition-title-bg}{definition-accent}{definition-title-fg},
        }{defn}
    \else
        \newtcbtheorem[number within=\ThemeTheoremCounterWithin]{mydefinition}{Definition}
        {
            theorem/variant={definition-body-bg}{definition-title-bg}{definition-accent}{definition-title-fg},
        }{defn}
    \fi

    \newtcbtheorem[use counter from=mydefinition]{mytheorem}{Theorem}
    {
        theorem/variant={theorem-body-bg}{theorem-title-bg}{theorem-accent}{theorem-title-fg},
    }{thm}

    \newtcbtheorem[use counter from=mydefinition]{mylemma}{Lemma}
    {
        theorem/variant={lemma-body-bg}{lemma-title-bg}{lemma-accent}{lemma-title-fg},
    }{lem}

    \newtcbtheorem[use counter from=mydefinition]{mycorollary}{Corollary}
    {
        theorem/variant={corollary-body-bg}{corollary-title-bg}{corollary-accent}{corollary-title-fg},
    }{cor}

    \newtcbtheorem[use counter from=mydefinition]{myproposition}{Proposition}
    {
        theorem/variant={proposition-body-bg}{proposition-title-bg}{proposition-accent}{proposition-title-fg},
    }{prop}

    \newtcbtheorem[use counter from=mydefinition]{myclaim}{Claim}
    {
        theorem/variant={claim-body-bg}{claim-title-bg}{claim-accent}{claim-title-fg},
    }{clm}

    \newtcbtheorem[use counter from=mydefinition]{myfact}{Fact}
    {
        theorem/variant={fact-body-bg}{fact-title-bg}{fact-accent}{fact-title-fg},
    }{fact}

    \newtcbtheorem[use counter from=mydefinition]{myassumption}{Assumption}
    {
        theorem/variant={assumption-body-bg}{assumption-title-bg}{assumption-accent}{assumption-title-fg},
    }{assump}
\fi

% -------------------- Study Callout Blocks --------------------
\NewDocumentEnvironment{insightbox}{O{Insight} +b}{
    \begin{tcolorbox}[callout/variant={insight-bg}{insight-accent}]
        {\textcolor{insight-label-fg}{\bfseries #1.}}\quad #2
}{
    \end{tcolorbox}
}

\NewDocumentEnvironment{pitfallbox}{O{Pitfall} +b}{
    \begin{tcolorbox}[callout/variant={pitfall-bg}{pitfall-accent}]
        {\textcolor{pitfall-label-fg}{\bfseries #1.}}\quad #2
}{
    \end{tcolorbox}
}

\NewDocumentEnvironment{intuitionbox}{O{Intuition} +b}{
    \begin{tcolorbox}[callout/variant={intuition-bg}{intuition-accent}]
        {\textcolor{intuition-label-fg}{\bfseries #1.}}\quad #2
}{
    \end{tcolorbox}
}

\NewDocumentEnvironment{summarybox}{O{Summary} +b}{
    \begin{tcolorbox}[callout/variant={summary-bg}{summary-accent}]
        {\textcolor{summary-label-fg}{\bfseries #1.}}\quad #2
}{
    \end{tcolorbox}
}

\NewDocumentEnvironment{questionbox}{O{Question} +b}{
    \begin{tcolorbox}[callout/variant={question-bg}{question-accent}]
        {\textcolor{question-label-fg}{\bfseries #1.}}\quad #2
}{
    \end{tcolorbox}
}

\NewDocumentCommand{\insight}{O{Insight} +m}{
    \begin{insightbox}[#1]
        #2
    \end{insightbox}
}

\NewDocumentCommand{\pitfall}{O{Pitfall} +m}{
    \begin{pitfallbox}[#1]
        #2
    \end{pitfallbox}
}

\NewDocumentCommand{\intuition}{O{Intuition} +m}{
    \begin{intuitionbox}[#1]
        #2
    \end{intuitionbox}
}

\NewDocumentCommand{\summary}{O{Summary} +m}{
    \begin{summarybox}[#1]
        #2
    \end{summarybox}
}

\NewDocumentCommand{\question}{O{Question} +m}{
    \begin{questionbox}[#1]
        #2
    \end{questionbox}
}

% -------------------- Theorem Commands --------------------
\ifEnablePlainAmsthmTheorem
    \NewDocumentCommand{\defn}{mm+m}{
        \begin{definition}\ThemeOptionalTheoremTitle{#1}
            \ThemeEmitPrefixedLabel{defn}{#2}
            #3
        \end{definition}
    }

    \NewDocumentCommand{\thm}{mm+m}{
        \begin{theorem}\ThemeOptionalTheoremTitle{#1}
            \ThemeEmitPrefixedLabel{thm}{#2}
            #3
        \end{theorem}
    }

    \NewDocumentCommand{\lem}{mm+m}{
        \begin{lemma}\ThemeOptionalTheoremTitle{#1}
            \ThemeEmitPrefixedLabel{lem}{#2}
            #3
        \end{lemma}
    }

    \NewDocumentCommand{\cor}{mm+m}{
        \begin{corollary}\ThemeOptionalTheoremTitle{#1}
            \ThemeEmitPrefixedLabel{cor}{#2}
            #3
        \end{corollary}
    }

    \NewDocumentCommand{\prop}{mm+m}{
        \begin{proposition}\ThemeOptionalTheoremTitle{#1}
            \ThemeEmitPrefixedLabel{prop}{#2}
            #3
        \end{proposition}
    }

    \NewDocumentCommand{\clm}{mm+m}{
        \begin{claim}\ThemeOptionalTheoremTitle{#1}
            \ThemeEmitPrefixedLabel{clm}{#2}
            #3
        \end{claim}
    }

    \newcommand{\ThemeFactCommand}[3]{
        \begin{fact}\ThemeOptionalTheoremTitle{#1}
            \ThemeEmitPrefixedLabel{fact}{#2}
            #3
        \end{fact}
    }

    \newcommand{\ThemeAssumptionCommand}[3]{
        \begin{assumption}\ThemeOptionalTheoremTitle{#1}
            \ThemeEmitPrefixedLabel{assump}{#2}
            #3
        \end{assumption}
    }

    \makeatletter
    \renewcommand{\fact}{\@ifnextchar\bgroup{\ThemeFactCommand}{\ThemeFactEnvBegin}}
    \renewcommand{\assumption}{\@ifnextchar\bgroup{\ThemeAssumptionCommand}{\ThemeAssumptionEnvBegin}}
    \makeatother
\else
    \NewDocumentCommand{\defn}{mm+m}{
        \ThemeRunTcbTheorem{mydefinition}{defn}{#1}{#2}{#3}
    }

    \NewDocumentCommand{\thm}{mm+m}{
        \ThemeRunTcbTheorem{mytheorem}{thm}{#1}{#2}{#3}
    }

    \NewDocumentCommand{\lem}{mm+m}{
        \ThemeRunTcbTheorem{mylemma}{lem}{#1}{#2}{#3}
    }

    \NewDocumentCommand{\cor}{mm+m}{
        \ThemeRunTcbTheorem{mycorollary}{cor}{#1}{#2}{#3}
    }

    \NewDocumentCommand{\prop}{mm+m}{
        \ThemeRunTcbTheorem{myproposition}{prop}{#1}{#2}{#3}
    }

    \NewDocumentCommand{\clm}{mm+m}{
        \ThemeRunTcbTheorem{myclaim*}{clm}{#1}{#2}{#3}
    }

    \NewDocumentCommand{\fact}{mm+m}{
        \ThemeRunTcbTheorem{myfact}{fact}{#1}{#2}{#3}
    }

    \NewDocumentCommand{\assumption}{mm+m}{
        \ThemeRunTcbTheorem{myassumption}{assump}{#1}{#2}{#3}
    }
\fi

\NewDocumentCommand{\lemp}{m+m+m}{
    \lem{#1}{}{#2}
    \pf{#3}
}

\NewDocumentCommand{\corp}{m+m+m}{
    \cor{#1}{}{#2}
    \pf{#3}
}

% Proof
\NewDocumentCommand{\pf}{+m}{
    \begin{proof}[\noindent\textbf{Proof.}]
        #1
    \end{proof}
}

\ifEnablePlainAmsthmTheorem
    % Example (plain mode)
    \newenvironment{example}{%
        \par
        \vspace{5pt}
        \noindent\textbf{Example. }\ignorespaces
    }{%
        \par
        \vspace{5pt}
    }

    \NewDocumentCommand{\ex}{+m}{
        \begin{example}
            #1
        \end{example}
    }

    % Remark (plain mode)
    \NewDocumentCommand{\rmk}{+m}{
        {\it #1}
    }

    \NewDocumentCommand{\rmkb}{+m}{
        \begin{remark}
            #1
        \end{remark}
    }

    \newenvironment{assump}{
        \par
        \vspace{5pt}
        \begingroup
        \noindent\textit{\textbf{Assumption. }}\ignorespaces
    }{
        \hbox{}\nobreak\hfill$\blacksquare$
        \par
        \endgroup
        \vspace{5pt}
    }

    % Note (plain mode)
    \NewDocumentEnvironment{note}{o g +b}{
        \par
        \vspace{5pt}
        \noindent
        \IfNoValueTF{#1}{
            \IfNoValueTF{#2}{
                \textbf{Note.}
            }{
                \textbf{Note (#2).}
            }
        }{
            \textbf{Note (#1).}
        }
        \par
        #3
    }{
        \par
        \vspace{5pt}
    }
\else
    % Example (styled mode)
    \newenvironment{example}{%
        \par
        \vspace{5pt}
        \begin{minipage}{\textwidth}
                \noindent\textbf{\color{example-label-fg}Example.}
                \tcolorbox[
                callout/variant={example-bg}{example-accent},
                callout/indented
            ]
    }{%
            \endtcolorbox
        \end{minipage}
        \vspace{5pt}
    }

    \NewDocumentCommand{\ex}{+m}{
        \begin{example}
            #1
        \end{example}
    }

    % Remark (styled mode)
    \NewDocumentCommand{\rmk}{+m}{
        {\it \color{remark-inline-fg}#1}
    }

    \newenvironment{styledremark}{
        \par
        \vspace{5pt}
        \begin{minipage}{\textwidth}
            {\par\noindent{\textbf{\color{remark-label-fg}Remark.}}}
            \tcolorbox[
                callout/variant={remark-bg}{remark-accent}
            ]
    }{
            \endtcolorbox
        \end{minipage}
        \vspace{5pt}
    }

    \NewDocumentCommand{\rmkb}{+m}{
        \begin{styledremark}
            #1
        \end{styledremark}
    }

    \newenvironment{assump}{
        {\noindent{\it \textbf{\color{assump-label-fg}Assumption.}}}
        \tcolorbox[
            callout/variant={assump-bg}{assump-accent},
            callout/indented
        ]
    }{
        \textcolor{assump-accent}{\hbox{}\nobreak\hfill$\blacksquare$}
        \endtcolorbox
    }

    % Note (styled mode)
    \NewDocumentEnvironment{note}{o g +b}{
        \IfNoValueTF{#1}{
            \IfNoValueTF{#2}{
                \begin{tcolorbox}[note/base]
            }{
                \begin{tcolorbox}[note/base,note/title={#2}]
            }
        }{
            \begin{tcolorbox}[note/base,note/title={#1}]
        }
        #3
    }{
        \end{tcolorbox}
    }
\fi
